% SOURCE docs/release/2026-10-29/source/docs/writing-repair/2026-10-29/src/45358a4637eac946.tex CHAR 2956
\chapter*{How to read the argument}
\addcontentsline{toc}{chapter}{How to read the argument}

The opening develops an annual deposit decision and four worked journeys.
Their cross-references distinguish exact supplied-law calculations, fitted
public and generated evidence, and unresolved identification. The longitudinal
and panel foundations are followed by the servicing and funding studies;
the observation foundations lead directly into the linkage bounds. The graph
journey points to the restricted temporal confidence construction and its
additional-information argument. The discovery foundations also develop the
completed search-repair comparison in Section~\ref{repair29:discovery}: a
computational stopping rule must not be mistaken for a scientifically
established relationship. The finite-sample comparison in
Section~
ef{closure30:sampling} then asks whether completing that search
produces reliable graphs and effect intervals. In the servicing study,
Sections~
ef{closure30:outcome-audit} and~
ef{rel73:partial-recovery}
distinguish buying more observations from recovering the intended historical
outcome. These distinctions connect computational and information choices
to the financial decision.

The evidence and method-choice part compares statistical targets before
selecting software. Chapter~\ref{scisyn:software} is the authoritative
software treatment: developed comparisons, the seventeen-area selection,
commercial alternatives and a target-weight illustration sit together.
Chapter~\ref{rewrite:assistance-tasks} brings formulation, measurement,
benchmark evidence and task selection into one assistance treatment.
The subsequent detailed parts supply source-specific derivations and
comparisons for readers examining a particular method or operating route.

The final investment part distinguishes the value of information from the
value of a common analytical product. Chapter~\ref{rewrite:investment-decision}
brings the annual benefit account, multi-year cost examples, institutional
variation and value-of-information calculation into one conditional decision.
Its comparison also prices the full effort per useful answer and the
waiting cost of batching in Section~
ef{closure:energy-study}. The
recommendation is tested against expert composition and existing
commercial capability. PDF bookmarks expose the detailed sections; chapter
and section references carry the reader between the opening question and
the supporting result.

Throughout the opening case, assignment of an offer, delivery of that offer,
acceptance and the resulting price are separate quantities. Balances and
contribution have different symbols and a stated one-year horizon. Other
journeys introduce their own units and decision dates explicitly. A number
from a constructed example is not a bank forecast. A method documented in an
inspected source is not thereby an executed integration. These distinctions
are essential to interpreting the proposal.

Research coverage retains the cutoff of 7 October 2026. Editions, acquisition
dates, inspected anchors and omissions remain in the source records. The
comparative recommendation does not claim that every literature or language
cell is complete. The unanswered questions that could change the first
investment are stated where that choice is made.


% SOURCE docs/release/2026-10-29/source/docs/causal_banking_release_candidate_2026-10-29.tex CHAR 1343

\begingroup\setlength{\parskip}{0pt}\tableofcontents\endgroup
\mainmatter
\part{The banking decision and the causal question}
\label{reading:part-decision}
A bank earns a return by changing a decision, not by observing a correlation.
The first task is to define the action, the alternative and the financial outcome.
Identification then determines what the data and assumptions can reveal;
estimation measures that quantity and its uncertainty.


